\section{Lei de Gauss (Magnetismo)}

Essa lei representa matematicamente o fato de não existirem monopolos magnéticos na natureza. A unidade básica do magnetismo é o dipolo magnético, que contém um polo norte (N) e um polo sul (S). A consequência disso para a configuração do campo magnético $\mathbf{B}$ é que não existe nenhum ponto no espaço que aja como fonte ou sorvedouro de $\mathbf{B}$. Pela discussão no apêndice B, isso pode ser representado matematicamente por
\begin{equation}
\nabla \cdot \mathbf{B} = 0.
\label{eq: dif_gauss_mag}
\end{equation}

Considere uma superfície gaussiana $S$ que delimita um volume $V$, integrando $\nabla \cdot \mathbf{B}$ em $V$, obtemos
\begin{equation}
\iiint_V \nabla \cdot \mathbf{B} \, dV = \oiint_S \mathbf{B} \cdot d\mathbf{S}.
\end{equation}

\noindent onde utilizamos o teorema da divergência (\ref{th: div}). Pela eq. \eqref{eq: dif_gauss_mag}, temos então 
\begin{equation}
\oiint_S \mathbf B \cdot d\mathbf{S} = 0.
\label{eq: int_gauss_mag}
\end{equation}

\noindent que no diz que o fluxo do campo magnético é nulo através de qualquer superfície gaussiana. 

Com efeito, se o campo magnético é gerado por uma distribuição espacial de dipolos magnéticos, o fluxo de $\mathbf{B}$ através de $S$ devido a cada dipolo é nulo, pois qualquer linha de força vinda de N que atravessa a superfície em uma orientação, deve atravessar de novo, mas com orientação oposta, para ir para S. Portanto, o fluxo total é zero. 
